\section{Conclusion and Future Work}\label{sec:conclusion}

We proposed a quantum interface that applies quantum minimum finding to the maximum flow network interdiction problem. It uses one qubit per arc, a Dicke-state preparation that enforces a cardinality budget, and a phase oracle that compares the follower's maximum flow with a threshold. With this interface, the quantum algorithm searches the same set of feasible plans as classical exhaustive enumeration and needs quadratically fewer accesses to the objective function, which is optimal in the query model (Proposition~\ref{prop:advantage}). Simulations on 30 instances agreed with the predicted square-root growth. We counted the gates of the interface components and ran amplitude amplification for a three-arc instance on an IBM processor, where the optimal plan appeared in 81.5\% of the shots against 25\% under uniform sampling. The speedup holds relative to exhaustive search, and the MILP remained much faster on every instance we could simulate.

Future work follows from these limitations. (i) A quantum speedup for branch-and-bound~\cite{montanaro2020} applied to the MILP would allow a comparison with the method used in practice. (ii) An explicit reversible maximum flow oracle with gate counts would replace the existence argument in Proposition~\ref{prop:advantage}. (iii) Resource estimates in the fault-tolerant regime would locate the crossover with classical enumeration and with the MILP. (iv) Efficient state preparations for general interdiction costs would extend the speedup beyond the cardinality variant. (v) A hardware run with the comparison against the noise model registered in advance would turn the exploratory analysis of Section~\ref{sec:results-hardware} into a test.
